\section{How a large language model works}
% Block 1 of PLAN.md, 20 min.  Message: an LLM is a very large fit that predicts
% the next token; each property of that fit has a practical consequence.

% ------------------------------------------------------------------------------
\begin{frame}{Text becomes tokens}{What the model actually sees}
  \vspace{2ex}
  \begin{center}
    \begin{tikzpicture}[tok/.style={draw=unibasPetrol,line width=0.6pt,
        rounded corners=1.5pt,inner sep=3pt,font=\small\ttfamily,
        minimum height=5mm}, node distance=1.2mm]
      \node[tok] (t1) {The};
      \node[tok,right=of t1] (t2) {\textvisiblespace Rabi};
      \node[tok,right=of t2] (t3) {\textvisiblespace frequency};
      \node[tok,right=of t3] (t4) {\textvisiblespace is};
      \node[tok,right=of t4] (t5) {\textvisiblespace 2};
      \node[tok,right=of t5] (t6) {$\pi$};
      \node[tok,right=of t6] (t7) {\textvisiblespace$\times$};
      \node[tok,right=of t7] (t8) {\textvisiblespace 5};
      \node[tok,right=of t8] (t9) {\textvisiblespace MHz};
      \onslide<2->{%
        \foreach \i/\n in {1/791,2/82341,3/11900,4/374,5/220,6/49345,7/25569,8/642,9/37594}
          \node[below=1.5mm of t\i,font=\scriptsize\ttfamily,text=unibasGrey] {\n};}
    \end{tikzpicture}
  \end{center}
  \vspace{1ex}

  \begin{itemize}
    \item<2-> Each token is an integer from a vocabulary of $\sim 10^5$
          \hfill{\scriptsize\color{unibasGrey}(numbers illustrative)}
    \item<3-> English: about $0.75$ words per token; code and \LaTeX{} need more
  \end{itemize}
\end{frame}

% ------------------------------------------------------------------------------
\begin{frame}{One job: predict the next token}
  \begin{equation*}
    p_\theta\left(x_{n+1} \,\middle|\, x_1, x_2, \ldots, x_n\right)
  \end{equation*}

  \begin{itemize}
    \item $x_1, \ldots, x_n$: the tokens so far, called the \emph{context}
    \item $p_\theta$: a probability distribution over the whole vocabulary
    \item<2-> $\theta$: the parameters of a neural network (a \emph{transformer}),
          $\sim 10^{11}$--$10^{12}$ of them
    \item<3-> Nothing else: no database lookup, no symbolic engine, no physics module
  \end{itemize}
\end{frame}

% ------------------------------------------------------------------------------
\begin{frame}{Generation is a loop}{Sample, append, repeat}
  \vspace{1ex}
  \begin{center}
    \begin{tikzpicture}[node distance=9mm]
      \node[box] (ctx) {\textbf{Context}\\ \ttfamily The Rabi frequency is};
      \node[box=figModel,right=of ctx] (net) {\textbf{Network} $p_\theta$};
      \node[box,right=of net] (smp) {\textbf{Sample}\\ \ttfamily\textvisiblespace 2};
      \draw[flow] (ctx) -- (net);
      \draw[flow] (net) -- (smp);
      \onslide<2->{\draw[flow] (smp.south) -- ++(0,-8mm) -| node[lbl,pos=0.25,below]
        {append the new token to the context} (ctx.south);}
    \end{tikzpicture}
  \end{center}
  \vspace{1ex}

  \begin{itemize}
    \item<2-> A long answer is thousands of passes through the loop
    \item<3-> In each pass, the context is all the model has to work with
    \item<4-> It cannot go back and revise a token, only write more
  \end{itemize}
\end{frame}

% ------------------------------------------------------------------------------
\begin{frame}{Training is a very large fit}{Pre-training}
  \begin{equation*}
    \mathcal{L}(\theta) = -\sum_{n} \ln p_\theta\left(x_{n+1} \,\middle|\, x_1, \ldots, x_n\right)
  \end{equation*}

  \begin{itemize}
    \item Minimise $\mathcal{L}$ over $\sim 10^{13}$ tokens: web, books, code, arXiv
    \item<2-> Think of a least-squares fit with $10^{12}$ free parameters
    \item<3-> Result: a compressed, lossy model of everything written, physics included
    \item<4-> It learns what text \emph{looks like}, not what is true
  \end{itemize}
\end{frame}

% ------------------------------------------------------------------------------
\begin{frame}{From text model to assistant}{Post-training}
  \begin{enumerate}
    \item \textbf{Pre-training}: continues any text; not yet an assistant
    \item<2-> \textbf{Instruction tuning}: imitates examples of good answers
    \item<3-> \textbf{Human feedback}: people rank answers; the model is trained
          towards the preferred ones (reinforcement learning from human feedback, RLHF)
    \item<4-> \textbf{Verifiable rewards}: maths with known answers, code that passes
          tests; this is where ``reasoning'' comes from
  \end{enumerate}

  \vspace{2ex}
  \onslide<5->{\centering\ckeybox[0.78\linewidth]{Rewarded for answers that
    \emph{look good to people} and \emph{pass checks}}\par}
\end{frame}

% ------------------------------------------------------------------------------
\begin{frame}{What follows from the training}
  \vspace{2ex}
  \renewcommand{\arraystretch}{1.6}
  \begin{tabularx}{\linewidth}{@{}p{0.36\linewidth}X@{}}
    \toprule
    \textbf{Property} & \textbf{Consequence for you} \\
    \midrule
    Trained to sound plausible
      & Right and wrong look the same; references get invented \\
    \reveal<2>{Trained to please (\emph{sycophancy})}
      & \reveal<2>{It agrees with you: ask it to find the flaw instead} \\
    \reveal<3>{Rewarded for passing checks}
      & \reveal<3>{It may weaken the check: read every change to a test} \\
    \bottomrule
  \end{tabularx}
\end{frame}

% ------------------------------------------------------------------------------
\begin{frame}{What follows from the architecture}
  \vspace{2ex}
  \renewcommand{\arraystretch}{1.6}
  \begin{tabularx}{\linewidth}{@{}p{0.36\linewidth}X@{}}
    \toprule
    \textbf{Property} & \textbf{Consequence for you} \\
    \midrule
    Random sampling
      & Same question, different answers: re-run and compare \\
    \reveal<2>{No memory between sessions}
      & \reveal<2>{Write things down: notes, instruction files, git history} \\
    \reveal<3>{Finite context}
      & \reveal<3>{What you put in decides what comes out} \\
    \bottomrule
  \end{tabularx}
\end{frame}

% ------------------------------------------------------------------------------
\begin{frame}[fragile]{Knowledge has a cutoff and a bias}{Julia is a minority language}
\begin{prompt}
Solve the master equation of a driven, decaying qubit with QuantumOptics.jl
\end{prompt}
\begin{julia}
using QuantumOptics
b  = SpinBasis(1//2)                   # |e⟩ = spin up, |g⟩ = spin down
Δ, Ω, Γ₁ = 0.0, 2π, 0.1                # rad/μs
H  = Δ/2*sigmaz(b) + Ω/2*sigmax(b)
ρ0 = dm(spindown(b))                   # start in |g⟩
ρt = mesolve(H, ρ0, 0:0.01:10, [sqrt(Γ₁)*sigmam(b)])   # QuTiP (Python) name
\end{julia}
\begin{onlyenv}<2->
\begin{julia}
tout, ρt = timeevolution.master(0:0.01:10, ρ0, H, [sqrt(Γ₁)*sigmam(b)])   # Julia
\end{julia}
\end{onlyenv}
  \begin{itemize}
    \item<3-> Far less Julia than Python in the training data: Python habits leak in
    \item<4-> Package interfaces change after the training cutoff
  \end{itemize}
\end{frame}
